\section[Cote 69 : les automorphismes du graphe des vingt-sept droites qui fixent un triangle]{Cote 69 : les automorphismes du graphe des vingt-sept droites qui fixent un triangle\protect\verifie{Folder69}}\label{sec:69}

La cote 69, \emph{Graphes cubiques : notes manuscrites (s.d.)}, cent dix-neuf pages, étudie le graphe $C$ des vingt-sept droites d'une surface cubique, deux droites étant liées quand elles se coupent, et son groupe d'automorphismes $W$, d'ordre $51\,840$. Aux pages~25 et~26, elle marque un triangle $t$ (trois droites deux à deux concourantes) et décrit $(C, t)$ par un ensemble $t$ à trois éléments, un plan $V$ sur $\mathbf F_2$ et un torseur sous une extension centrale $\widetilde\Psi$ d'ordre $32$, produit de Baer de trois copies de $Q_8$ (\lot{69}{2}{25}, \lot{69}{2}{26}). La page~25 donne le groupe $W'_t$ des automorphismes qui fixent $t$ point par point comme extension de $\mathfrak S_4$ par $\mathbf F_2^3$, et le corollaire de la page~26 en tire l'ordre du groupe $W_t$ des automorphismes de $(C, t)$.

\noindent\emph{Conventions.} Les sommets sont $a_i$, $b_i$ ($0 \leq i \leq 5$) et $c_{ij}$ ($i < j$) ; $a_i$ et $b_j$ sont liés si $i \neq j$, $a_i$ et $b_i$ sont liés à $c_{jk}$ si $i \in \{j, k\}$, $c_{ij}$ et $c_{kl}$ sont liés si $\{i, j\} \cap \{k, l\} = \emptyset$, et il n'y a pas d'autre arête. C'est le modèle classique des vingt-sept droites ; chaque sommet a dix voisins. On prend $t = \{a_0, b_1, c_{01}\}$.

\begin{theoreme}[Grothendieck, page 26]
Le groupe $W_t$ des automorphismes de $C$ qui conservent le triangle $t$ est d'ordre $32 \cdot 36 = 1152$.
\end{theoreme}

\noindent L'énoncé tient : voir \texttt{card\_Wt} et \texttt{card\_aut} dans \href{\leanblob/Folder69.lean}{\texttt{Folder69.lean}}. Le sous-groupe $W'_t$ est d'ordre $192 = 8 \cdot 24$, comme le donne la suite exacte de la page~25 (\texttt{card\_pointwise}). La vérification est un calcul fini : la base $a_0, b_1, a_1, a_2, a_3, a_4$ a des orbites de tailles $3, 2, 8, 4, 3, 2$ ; chaque orbite est contenue dans l'ensemble des sommets qui ont les mêmes incidences avec les points déjà fixés, et parcourue par des automorphismes explicites que le noyau vérifie ; un automorphisme qui fixe la base fixe tout, car la base sépare les sommets. La cote~74 retrouve $1152$ par un autre chemin, les $96$ autotopies de la table de Klein (§\,\ref{sec:74}).

Ne sont pas démontrés ici : la description de $(C, t)$ par $(t, V, \widetilde P)$ et l'équivalence de catégories de la page~26, l'identification de $\widetilde\Psi$ au produit de Baer des relevés de $Q_8$, la suite exacte $1 \to \widetilde\Psi \to W_t \to \mathfrak S_t \times \mathfrak S_{V^*} \to 1$, et le rapprochement de $W_t$ avec $W(F_4)$.
